\chapter{PyTorch and the Machine Underneath}\label{ch:pytorch}

Every model in this book is a PyTorch program, and every PyTorch program is a stream of \emph{kernels}
launched onto an accelerator. This chapter covers both levels. The first half is a working introduction
to PyTorch: tensors, automatic differentiation, modules, data loaders, the training loop, checkpoints, and
running on one or several GPUs. A reader fluent in PyTorch can skim it, but should read the Systems Lens
boxes, which connect each API to what happens on the hardware. The second half opens the box: how a GPU
is organised, why most kernels are limited by memory bandwidth rather than arithmetic, how the
\emph{roofline model} predicts which, how PyTorch's asynchronous execution model hides (or fails to hide)
latency, and how to measure all of this honestly. These are the tools we use to reason about every
optimisation in later chapters.

\begin{figure}[t]
\centering
\begin{tikzpicture}[
  stage/.style={draw=sysblue,rounded corners=3pt,fill=sysbluebg,minimum width=2.05cm,
    minimum height=1.0cm,align=center,font=\footnotesize},
  hw/.style={stage,fill=napkinbg,draw=napkin},
  num/.style={circle,fill=sysblue,text=white,font=\scriptsize\bfseries,inner sep=1.2pt},
  arr/.style={-{Stealth[length=2.2mm]},thick,gray!70}]
\node[stage] (a) {Tensors \&\\dtypes};
\node[stage,right=0.35cm of a] (b) {Autograd \&\\graphs};
\node[stage,right=0.35cm of b] (c) {Modules,\\loaders, loop};
\node[stage,right=0.35cm of c] (d) {GPUs \&\\DDP};
\node[hw,below=0.9cm of a] (e) {GPU\\architecture};
\node[hw,right=0.35cm of e] (f) {Roofline\\model};
\node[hw,right=0.35cm of f] (g) {Async execution\\\& compile};
\node[hw,right=0.35cm of g] (h) {Measuring \&\\profiling};
\foreach \n/\k in {a/1,b/2,c/3,d/4,e/5,f/6,g/7,h/8} \node[num] at (\n.north west) {\k};
\draw[arr] (a)--(b); \draw[arr] (b)--(c); \draw[arr] (c)--(d);
\draw[arr] (e)--(f); \draw[arr] (f)--(g); \draw[arr] (g)--(h);
\draw[arr] (d.south) |- ($(d.south)!0.5!(h.north)$) -| (e.north);
\node[font=\scriptsize,sysblue,anchor=east] at ($(a.west)+(-0.1,0)$) {PyTorch};
\node[font=\scriptsize,napkin,anchor=east] at ($(e.west)+(-0.1,0)$) {hardware};
\end{tikzpicture}
\caption{Roadmap of \cref{ch:pytorch}: four stages of PyTorch fundamentals followed by four stages that
explain what those programs do on the hardware.}
\label{fig:pytorch:roadmap}
\end{figure}

% =====================================================================
\section{What is PyTorch?}\label{sec:pytorch:what}

PyTorch is an open-source Python library for deep learning, originally released by Meta AI in 2016 and now
governed by the PyTorch Foundation. It has been the most widely used framework in machine-learning research
for years, and it is the substrate of nearly every open LLM codebase, from Hugging Face Transformers to
vLLM, SGLang, Megatron-LM and veRL. Its popularity comes from a design that feels like ordinary Python: you
write imperative code, run it line by line, inspect any value, and get gradients for free.

\subsection{The three core components}
It is helpful to see PyTorch as three libraries in one (\cref{fig:pytorch:stack}):
\begin{enumerate}
\item A \textbf{tensor library}: an array library in the spirit of NumPy, whose arrays can live on CPUs,
  NVIDIA and AMD GPUs, or Apple silicon, with a vast set of operations implemented as optimised kernels.
\item An \textbf{automatic differentiation engine} (\emph{autograd}): it records the operations applied to
  tensors and computes gradients by traversing that record backwards.
\item A \textbf{deep-learning library}: modules (layers), loss functions, optimisers, data loading and
  distributed training utilities built on top of the first two.
\end{enumerate}

\begin{figure}[t]
\centering
\begin{tikzpicture}[font=\footnotesize,
  lay/.style={draw,minimum width=9.6cm,minimum height=0.62cm,align=center},
  side/.style={font=\scriptsize,align=left,anchor=west,text width=3.6cm}]
\node[lay,fill=keygreenbg,draw=keygreen] (py) at (0,3.6) {\textbf{Your model code}: \texttt{nn.Module}, training loop};
\node[lay,fill=keygreenbg,draw=keygreen] (lib) at (0,2.9) {\texttt{torch.nn}, \texttt{torch.optim}, \texttt{torch.utils.data}, \texttt{torch.distributed}};
\node[lay,fill=sysbluebg,draw=sysblue] (ag) at (0,2.2) {autograd: records ops, runs backward};
\node[lay,fill=sysbluebg,draw=sysblue] (disp) at (0,1.5) {dispatcher: picks a kernel by device, dtype, layout (ATen)};
\node[lay,fill=napkinbg,draw=napkin] (k) at (0,0.8) {kernels: cuBLAS, cuDNN, FlashAttention, Triton, oneDNN, \ldots};
\node[lay,fill=napkinbg,draw=napkin] (drv) at (0,0.1) {CUDA runtime \& driver $\to$ GPU (streams, allocator)};
\node[side] at (5.0,3.25) {Python, eager};
\node[side] at (5.0,1.85) {C++ core};
\node[side] at (5.0,0.45) {vendor libraries / generated code};
\draw[decorate,decoration={brace,amplitude=4pt},thick] (4.9,3.95) -- (4.9,2.6);
\draw[decorate,decoration={brace,amplitude=4pt},thick] (4.9,2.5) -- (4.9,1.2);
\draw[decorate,decoration={brace,amplitude=4pt},thick] (4.9,1.1) -- (4.9,-0.2);
\draw[-{Stealth},thick,labpurple] (-5.1,3.6) -- (-5.1,0.1);
\node[font=\scriptsize,align=right,text=labpurple,anchor=east] at (-5.25,1.85) {one\\Python call\\becomes one\\or more kernel\\launches};
\end{tikzpicture}
\caption{The PyTorch stack. A single Python-level operation passes through autograd and the dispatcher
before a device kernel is enqueued. Each layer adds a few microseconds of overhead, which is why tiny
operations are slow in eager mode and why kernel fusion and compilation matter (\cref{sec:pytorch:async}).}
\label{fig:pytorch:stack}
\end{figure}

\subsection{Deep learning in one paragraph}
Artificial intelligence aims at systems that perform tasks associated with human intelligence; machine
learning is the subfield that achieves this by learning from data; deep learning is the part of machine
learning that uses neural networks with many layers, which learn both the features and the decision rule
from raw inputs. Training a deep network follows a fixed recipe: a model with parameters $\theta$ maps inputs
to predictions; a loss function measures how wrong the predictions are on labelled training data; gradients
of the loss with respect to $\theta$ are computed by backpropagation; and an optimiser nudges $\theta$ against the
gradient. After training, the model is used for \emph{inference} on new inputs. Everything in this book,
including LLM pretraining and RL, is an instance of this recipe with a particular model, loss and data stream.

\subsection{Installing PyTorch}
PyTorch is installed with \code{pip} or \code{conda}. CPU-only wheels work everywhere; GPU builds must match
the CUDA (or ROCm) version supported by the installed driver, and the selector on \code{pytorch.org} prints the
right command. Two lines confirm what is available:

\begin{lstlisting}[caption={Checking the installation and available accelerators.},label={lst:pytorch:check}]
import torch
print(torch.__version__)
print("CUDA:", torch.cuda.is_available(),
      "| MPS:", torch.backends.mps.is_available())   # Apple silicon GPU
if torch.cuda.is_available():
    p = torch.cuda.get_device_properties(0)
    print(p.name, f"{p.total_memory/2**30:.0f} GiB,", p.multi_processor_count, "SMs")
\end{lstlisting}

Apple-silicon Macs expose their GPU through the \code{mps} backend, which is convenient for the small
experiments in this book but lacks some kernels and all of the CUDA-specific tooling; we write code that
selects a device with \code{llmsys.utils.get\_device()} and runs on whatever is present. \Cref{app:setup}
gives a complete environment recipe, including cloud GPUs.

% =====================================================================
\section{Tensors}\label{sec:pytorch:tensors}

\subsection{Scalars, vectors, matrices and beyond}
A tensor is an $n$-dimensional array with a single element type. Its \emph{rank} (or order) is the number
of dimensions: a scalar has rank 0, a vector rank 1, a matrix rank 2, and LLM activations are usually rank 3
with shape \texttt{(batch, sequence, features)} or rank 4 inside attention,
\texttt{(batch, heads, sequence, head\_dim)}. Creating tensors is straightforward:

\begin{lstlisting}[caption={Creating tensors of different ranks.},label={lst:pytorch:create}]
s = torch.tensor(3)                      # rank 0, shape ()
v = torch.tensor([1.0, 2.0, 3.0])        # rank 1, shape (3,)
m = torch.tensor([[1, 2], [3, 4]])       # rank 2, shape (2, 2), int64
x = torch.zeros(8, 1024, 768)            # rank 3: (batch, seq, d_model)
print(x.ndim, x.shape, x.dtype, x.device, x.numel())
\end{lstlisting}

\subsection{Tensor data types}
PyTorch infers the data type (\emph{dtype}) from the input: Python integers become \code{torch.int64}, Python
floats become \code{torch.float32}. Token ids are integers (\code{int64} by default, since embedding lookups
expect it); everything that is learned or computed is floating point. Converting is explicit:
\code{x.to(torch.bfloat16)} or \code{x.float()}. The default \code{float32} is a sensible balance of precision
and cost on CPUs, but on modern accelerators most LLM arithmetic runs in 16-bit or even 8-bit formats, which
we examine in \cref{sec:pytorch:formats}.

\subsection{Common tensor operations}
A handful of operations recur throughout the book:

\begin{lstlisting}[caption={Operations used constantly in LLM code.},label={lst:pytorch:ops}]
a = torch.arange(6.0).reshape(2, 3)     # [[0,1,2],[3,4,5]]
a.view(3, 2)                            # reinterpret same memory (needs contiguity)
a.reshape(3, 2)                         # view if possible, else copy
a.T                                     # transpose of a 2-D tensor (a view)
a.transpose(0, 1)                       # swap any two dims (a view)
a @ a.T                                 # matrix multiplication, (2,3)@(3,2) -> (2,2)
q = torch.randn(4, 12, 128, 64)         # (batch, heads, seq, head_dim)
scores = q @ q.transpose(-2, -1)        # batched matmul -> (4, 12, 128, 128)
x = torch.randn(2, 5, 8)
x.mean(dim=-1, keepdim=True)            # reductions along a dimension
torch.softmax(x, dim=-1)                # normalise the last dimension
torch.cat([x, x], dim=1)                # (2, 10, 8)
x[:, -1, :]                             # last position of every sequence, (2, 8)
\end{lstlisting}

\begin{systemslens}[Shapes, strides and the cost of a ``free'' transpose]
A tensor is a pointer into a block of memory (its \emph{storage}) plus a shape and a \emph{stride} per
dimension: the number of elements to skip to move one step along that dimension. A contiguous
$2\times3$ tensor has strides $(3,1)$; its transpose is the \emph{same storage} with shape $(3,2)$ and strides
$(1,3)$. That is why \code{.T}, \code{.transpose}, slicing and \code{.view} are free: they allocate nothing and
launch no kernel. The bill arrives later. A kernel that reads a non-contiguous tensor makes strided,
uncoalesced memory accesses, and some operations (\code{.view} after a transpose, many custom kernels) refuse
non-contiguous inputs, forcing a \code{.contiguous()} copy. In attention code, the
\code{transpose(1, 2)} between \texttt{(B,T,H,D)} and \texttt{(B,H,T,D)} layouts is the classic source of
hidden copies. Lab~2 asks you to detect views versus copies by comparing storage pointers.
\end{systemslens}

\subsection{Number formats}\label{sec:pytorch:formats}
Floating-point formats trade \emph{range} (exponent bits) against \emph{precision} (mantissa bits)
(\cref{fig:pytorch:formats}). The formats that matter for LLMs are:
\begin{itemize}
\item \textbf{fp32} (1 sign, 8 exponent, 23 mantissa bits): the reference format; master weights and
  optimiser states are usually kept in fp32.
\item \textbf{tf32}: not a storage format but a tensor-core \emph{compute mode} that reads fp32 inputs and
  multiplies them with a 10-bit mantissa, enabled in PyTorch by setting
  \code{torch.backends.cuda.matmul.allow\_tf32} to \code{True}.
\item \textbf{fp16} (1/5/10): good precision but a narrow range (maximum 65{,}504), so gradients can
  underflow or overflow; training in fp16 needs \emph{loss scaling} (\cref{ch:pretrain}).
\item \textbf{bf16} (1/8/7): the same range as fp32 with much less precision (machine epsilon $2^{-7}$). Its
  range makes it the default for LLM training and inference.
\item \textbf{fp8} in two flavours, E4M3 (max 448, more precision) and E5M2 (max 57{,}344, more range),
  supported by Hopper and later GPUs for matmuls with per-tensor or per-block scaling factors
  (\cref{ch:efficient}).
\end{itemize}

\begin{figure}[t]
\centering
\begin{tikzpicture}[font=\scriptsize,x=0.25cm,y=0.5cm]
\newcommand{\fmt}[5]{% y, name, exp bits, mant bits, note
  \node[anchor=east,font=\footnotesize\bfseries] at (-0.4,#1) {#2};
  \fill[warnred!60] (0,#1-0.35) rectangle (1,#1+0.35);
  \fill[sysblue!45] (1,#1-0.35) rectangle ({1+#3},#1+0.35);
  \fill[keygreen!45] ({1+#3},#1-0.35) rectangle ({1+#3+#4},#1+0.35);
  \draw (0,#1-0.35) rectangle ({1+#3+#4},#1+0.35);
  \foreach \b in {1,...,\numexpr#3+#4\relax} \draw[gray!50] (\b,#1-0.35) -- (\b,#1+0.35);
  \node[anchor=west] at (33.5,#1) {#5};}
\fmt{0}{fp32}{8}{23}{$\pm3.4\times10^{38}$, eps $1.2\times10^{-7}$}
\fmt{-1.2}{tf32}{8}{10}{matmul mode only}
\fmt{-2.4}{bf16}{8}{7}{fp32 range, eps $7.8\times10^{-3}$}
\fmt{-3.6}{fp16}{5}{10}{max 65{,}504, eps $9.8\times10^{-4}$}
\fmt{-4.8}{fp8 E4M3}{4}{3}{max 448}
\fmt{-6.0}{fp8 E5M2}{5}{2}{max 57{,}344}
\node[anchor=west] at (0,1.0) {\textcolor{warnred}{\rule{6pt}{6pt}} sign \quad \textcolor{sysblue}{\rule{6pt}{6pt}} exponent (range) \quad \textcolor{keygreen}{\rule{6pt}{6pt}} mantissa (precision)};
\end{tikzpicture}
\caption{Bit layouts of the floating-point formats used for LLMs. bf16 keeps fp32's eight exponent bits and
sacrifices mantissa; fp16 does the opposite. Fewer bits mean fewer bytes to move and more tensor-core
throughput: on an H100, dense bf16 matmul peaks at about 989\,TFLOP/s and fp8 at about twice that.}
\label{fig:pytorch:formats}
\end{figure}

\begin{napkinmath}[What a dtype buys you]
Halving the bytes per element halves memory capacity and traffic for every tensor stored in that format,
and on tensor cores roughly doubles peak matmul throughput (H100 SXM, dense: tf32 $\approx$495, bf16 $\approx$989,
fp8 $\approx$1{,}979\,TFLOP/s). A 7B model's weights take 28\,GB in fp32, 14\,GB in bf16, 7\,GB in fp8, and
3.5\,GB in int4. Precision is the cost: bf16's relative spacing of $2^{-7}\approx0.8\%$ means that adding a
small update to a large weight can round to no change at all. This is why optimisers keep fp32 master
weights even when the forward pass runs in bf16 (\cref{ch:pretrain}).
\end{napkinmath}

% =====================================================================
\section{Models as computation graphs}\label{sec:pytorch:graphs}

A neural network is a composition of operations, and it is useful to draw it as a directed acyclic graph
whose nodes are operations and whose edges carry tensors: the \emph{computation graph}. Take logistic
regression on one example, the smallest model with every ingredient of a deep network:
\begin{equation}
z = \vect{w}^\top\vect{x} + b,\qquad a = \sigma(z) = \frac{1}{1+e^{-z}},\qquad
\mathcal{L} = -\bigl[y\log a + (1-y)\log(1-a)\bigr].
\end{equation}

\begin{lstlisting}[caption={Logistic regression forward pass; autograd records the graph as it runs.},label={lst:pytorch:logreg}]
import torch.nn.functional as F
x = torch.tensor([1.5, -0.5])       # input features
y = torch.tensor(1.0)               # true label
w = torch.tensor([0.4, 0.9], requires_grad=True)
b = torch.tensor(0.1, requires_grad=True)
z = x @ w + b                       # weighted sum
a = torch.sigmoid(z)                # predicted probability
loss = F.binary_cross_entropy(a, y)
print(loss.grad_fn)                 # <BinaryCrossEntropyBackward0 ...>
\end{lstlisting}

PyTorch builds this graph \emph{dynamically} (``define-by-run''): each operation on a tensor with
\code{requires\_grad=True} appends a node, recorded in the result's \code{grad\_fn}, that knows how to compute
its local derivative. \Cref{fig:pytorch:graph} shows the graph for \cref{lst:pytorch:logreg} together with
the backward pass that the next section runs over it.

\begin{figure}[t]
\centering
\begin{tikzpicture}[font=\footnotesize,
  op/.style={circle,draw=sysblue,fill=sysbluebg,minimum size=0.85cm,inner sep=1pt},
  leaf/.style={draw,rounded corners=2pt,fill=keygreenbg,draw=keygreen,minimum height=0.55cm,inner sep=3pt},
  dat/.style={draw,rounded corners=2pt,fill=gray!10,draw=gray,minimum height=0.55cm,inner sep=3pt},
  g/.style={draw,rounded corners=2pt,fill=warnredbg,draw=warnred,minimum height=0.55cm,inner sep=3pt,font=\scriptsize},
  f/.style={-{Stealth[length=2mm]},thick},
  bk/.style={-{Stealth[length=2mm]},thick,warnred,dashed}]
\node[dat] (x) at (0,1.0) {$\vect{x}$};
\node[leaf] (w) at (0,0) {$\vect{w}$};
\node[op] (mm) at (1.9,0.5) {$\cdot$};
\node[leaf] (b) at (1.9,1.6) {$b$};
\node[op] (add) at (3.8,0.5) {$+$};
\node[op] (sig) at (5.7,0.5) {$\sigma$};
\node[dat] (y) at (7.6,1.6) {$y$};
\node[op] (bce) at (7.6,0.5) {BCE};
\node[draw,rounded corners=2pt,fill=warnredbg,draw=warnred,inner sep=3pt] (L) at (9.4,0.5) {$\mathcal{L}$};
\draw[f] (x) -- (mm); \draw[f] (w) -- (mm);
\draw[f] (mm) -- node[above] {$u$} (add); \draw[f] (b) -- (add);
\draw[f] (add) -- node[above] {$z$} (sig); \draw[f] (sig) -- node[above] {$a$} (bce);
\draw[f] (y) -- (bce); \draw[f] (bce) -- (L);
% backward lane
\node[g] (gL) at (9.4,-1.3) {$1$};
\node[g] (ga) at (7.6,-1.3) {$\partial\mathcal{L}/\partial a$};
\node[g] (gz) at (3.8,-1.3) {$\partial\mathcal{L}/\partial z=a-y$};
\node[g] (gw) at (0.6,-2.5) {$\partial\mathcal{L}/\partial \vect{w}=(a-y)\,\vect{x}$};
\node[g] (gb) at (4.6,-2.5) {$\partial\mathcal{L}/\partial b=a-y$};
\draw[bk] (gL) -- (ga); \draw[bk] (ga) -- node[above,font=\scriptsize] {$\times\,a(1-a)$} (gz);
\draw[bk] (gz.south west) -- (gw.north east); \draw[bk] (gz) -- (gb);
\draw[gray!50,dotted] (L) -- (gL); \draw[gray!50,dotted] (sig) -- (ga.north west); \draw[gray!50,dotted] (add) -- (gz);
\draw[gray!50,dotted] (w) -- (gw); \node[font=\scriptsize,anchor=west,warnred] at (7.6,-2.5) {backward pass};
\node[font=\scriptsize,anchor=west] at (8.2,1.6) {forward pass};
\end{tikzpicture}
\caption{Computation graph for logistic regression. Solid arrows: forward pass (tensors flow left to right).
Dashed arrows: the backward pass multiplies local derivatives from the loss back to the leaves that
require gradients ($\vect{w}$ and $b$, green). For sigmoid followed by BCE, the product simplifies to
$\partial\mathcal{L}/\partial z = a-y$.}
\label{fig:pytorch:graph}
\end{figure}

% =====================================================================
\section{Automatic differentiation}\label{sec:pytorch:autograd}

Training needs the gradient of the loss with respect to every parameter. By the chain rule, the gradient
of $\mathcal{L}$ with respect to $\vect{w}$ is a product of local derivatives along the path from $\vect{w}$
to $\mathcal{L}$. \emph{Reverse-mode} automatic differentiation evaluates that product from the loss
backwards, which costs about the same as one more forward pass no matter how many parameters there are.
This is the right mode for a scalar loss and millions of parameters.

PyTorch exposes reverse mode in two ways. \code{torch.autograd.grad} returns gradients of an output with
respect to chosen inputs; \code{.backward()} accumulates gradients into the \code{.grad} field of every leaf
that requires them:

\begin{lstlisting}[caption={Two ways to get gradients from the graph of Listing~\ref{lst:pytorch:logreg}.},label={lst:pytorch:grad}]
from torch.autograd import grad
gw = grad(loss, w, retain_graph=True)     # keep the graph for a second call
gb = grad(loss, b, retain_graph=True)
loss.backward()                           # fills w.grad and b.grad (and frees the graph)
print(w.grad, b.grad)                     # equal to gw[0], gb[0]
print((a - y).detach() * x)               # the analytic gradient for w
\end{lstlisting}

Three details cause most autograd bugs. First, gradients \emph{accumulate}: \code{.backward()} adds to
\code{.grad}, so training loops must zero gradients every step (or rely on accumulation deliberately for
gradient accumulation, \cref{ch:pretrain}). Second, the graph is freed after one backward pass unless
\code{retain\_graph=True}. Third, operations under \code{torch.no\_grad()} or on \code{.detach()}ed tensors are not
recorded, which is what we want during evaluation and generation.

\begin{systemslens}[Why backward costs twice the forward, and where activation memory comes from]
For a linear layer $\mat{Y}=\mat{X}\mat{W}$ with $\mat{X}\in\R^{m\times k}$ and $\mat{W}\in\R^{k\times n}$, the
forward pass costs $2mkn$ FLOPs. Backward needs two products of the same size:
$\partial\mathcal{L}/\partial\mat{X}=\mat{G}\mat{W}^\top$ and $\partial\mathcal{L}/\partial\mat{W}=\mat{X}^\top\mat{G}$,
where $\mat{G}=\partial\mathcal{L}/\partial\mat{Y}$, for $4mkn$ FLOPs. Hence forward plus backward
$\approx 3\times$ forward, and with $2N$ FLOPs per token forward we get the $6N$ FLOPs per training token used
throughout the book. The second product needs $\mat{X}$, the layer's \emph{input}, so autograd must keep
every such input alive from the forward pass until its backward step: these \emph{saved activations} are why
training memory grows with batch size $\times$ sequence length $\times$ depth, and why activation
checkpointing (recomputing them during backward, \cref{ch:pretrain}) trades FLOPs for memory.
\end{systemslens}

% =====================================================================
\section{Implementing multilayer neural networks}\label{sec:pytorch:mlp}

Real models are built by subclassing \code{torch.nn.Module}. A module registers its parameters (as
\code{nn.Parameter} attributes or submodules) in \code{\_\_init\_\_} and defines the computation in
\code{forward}. Calling the module runs \code{forward} and, via hooks, integrates with autograd, device
movement, \code{state\_dict}s and more.

\begin{lstlisting}[caption={A multilayer perceptron as an \texttt{nn.Module}.},label={lst:pytorch:mlp}]
import torch.nn as nn

class MLP(nn.Module):
    def __init__(self, in_dim, out_dim, hidden=(30, 20)):
        super().__init__()
        layers, prev = [], in_dim
        for h in hidden:
            layers += [nn.Linear(prev, h), nn.ReLU()]
            prev = h
        layers.append(nn.Linear(prev, out_dim))   # outputs logits
        self.net = nn.Sequential(*layers)

    def forward(self, x):
        return self.net(x)

torch.manual_seed(123)                    # reproducible random initialisation
model = MLP(50, 3)
print(sum(p.numel() for p in model.parameters() if p.requires_grad))   # 2213
\end{lstlisting}

The count is $(50\cdot30+30)+(30\cdot20+20)+(20\cdot3+3)=2{,}213$: each \code{nn.Linear(in, out)} has an
$\text{out}\times\text{in}$ weight matrix and an \code{out}-dimensional bias. The weights are initialised randomly
(by default from a scaled uniform distribution), so fixing the seed makes runs reproducible.

Two habits are worth forming early. First, the network returns \emph{logits}, unnormalised scores, not
probabilities: PyTorch's \code{cross\_entropy} applies a numerically stable log-softmax internally, and
applying softmax yourself first is both wasteful and less accurate. Call \code{torch.softmax(logits, dim=-1)}
only when you actually need probabilities, for instance for sampling. Second, wrap inference in
\code{torch.no\_grad()} (or \code{torch.inference\_mode()}), so that no graph is recorded and no activations are kept:

\begin{lstlisting}[caption={Inference without building a graph.},label={lst:pytorch:nograd}]
X = torch.rand(4, 50)
with torch.no_grad():
    logits = model(X)
    probs = torch.softmax(logits, dim=-1)   # rows sum to 1
\end{lstlisting}

% =====================================================================
\section{Setting up efficient data loaders}\label{sec:pytorch:loaders}

PyTorch separates \emph{what} the data are from \emph{how} they are batched. A \code{Dataset} implements
\code{\_\_len\_\_} and \code{\_\_getitem\_\_(i)}, returning one example. A \code{DataLoader} draws indices (shuffled
or not), fetches examples (optionally in worker processes), collates them into batch tensors, and yields
them:

\begin{lstlisting}[caption={A map-style dataset and its loader.},label={lst:pytorch:dataset}]
from torch.utils.data import Dataset, DataLoader

class ToyDataset(Dataset):
    def __init__(self, X, y):
        self.features, self.labels = X, y
    def __getitem__(self, i):
        return self.features[i], self.labels[i]
    def __len__(self):
        return self.labels.shape[0]

train_loader = DataLoader(ToyDataset(X_train, y_train), batch_size=2,
                          shuffle=True,      # new order every epoch
                          drop_last=True,    # avoid a tiny final batch
                          num_workers=0)     # 0 = load in the main process
for epoch in range(2):
    for xb, yb in train_loader:
        ...
\end{lstlisting}

\code{shuffle=True} matters for SGD, which assumes batches are roughly independent samples; \code{drop\_last=True}
avoids a final, much smaller batch whose noisy gradient can destabilise training (and which would also break
fixed-shape compilation). \code{num\_workers} controls parallel loading.

\begin{figure}[t]
\centering
\resizebox{\linewidth}{!}{%
\begin{tikzpicture}[font=\footnotesize,
  w/.style={draw=keygreen,fill=keygreenbg,rounded corners=2pt,minimum width=1.9cm,minimum height=0.55cm},
  q/.style={draw=gray,fill=gray!10,minimum width=0.45cm,minimum height=0.45cm},
  g/.style={draw=sysblue,fill=sysbluebg,rounded corners=2pt,minimum width=2.0cm,minimum height=1.6cm,align=center},
  arr/.style={-{Stealth[length=1.8mm]},thick,gray!70}]
\node[draw=napkin,fill=napkinbg,rounded corners=2pt,minimum height=1.9cm,minimum width=1.6cm,align=center] (disk) at (0,0) {storage\\(files,\\memmap)};
\foreach \i in {0,1,2} \node[w] (w\i) at (3.0,0.75-\i*0.75) {worker \the\numexpr\i+1\relax};
\foreach \i in {0,1,2,3} \node[q] (q\i) at (5.0+\i*0.5,0) {};
\node[font=\scriptsize] at (5.75,0.5) {prefetch queue};
\node[draw=labpurple,fill=labpurplebg,rounded corners=2pt,minimum width=1.4cm,minimum height=0.9cm,align=center] (pin) at (9.0,0) {pinned\\host buf};
\node[g] (gpu) at (12.3,0) {GPU\\compute};
\foreach \i in {0,1,2} {\draw[arr] (disk.east) -- (w\i.west); \draw[arr] (w\i.east) -- (q0.west);}
\draw[arr] (q3.east) -- (pin.west) node[midway,above=2pt,font=\scriptsize] {collate};
\draw[arr,labpurple] (pin.east) -- (gpu.west) node[midway,above=2pt,font=\scriptsize] {async H2D};
\node[font=\scriptsize,align=center] at (3.0,-1.45) {read, decode,\\tokenise, augment};
\node[font=\scriptsize,align=center] at (12.3,-1.2) {must never wait};
\end{tikzpicture}}
\caption{A data pipeline. Worker processes produce batches ahead of time into a bounded queue; with
\code{pin\_memory=True} batches land in page-locked host memory so the host-to-device copy can be asynchronous
and overlap with compute.}
\label{fig:pytorch:loader}
\end{figure}

\begin{systemslens}[Feeding the accelerator]
A DataLoader with \code{num\_workers=k} forks $k$ processes (side-stepping Python's global interpreter lock),
each producing whole batches, and keeps \code{prefetch\_factor} batches per worker in flight
(\cref{fig:pytorch:loader}). The goal is simple: the GPU should never wait for input. For LLM pretraining the
per-batch CPU work is tiny if, and only if, tokenisation happened offline: a batch is then just a slice of a
memory-mapped array of token ids (\code{llmsys.data.TokenMemmapDataset}, \cref{ch:text}). Two further flags
matter on GPUs: \code{pin\_memory=True} allocates batches in page-locked memory, which the DMA engine can copy
from directly, and \code{x.to(device, non\_blocking=True)} then makes the copy asynchronous so it overlaps with
the previous step's compute. In Jupyter notebooks, multi-process loading is unreliable; use
\code{num\_workers=0} there and in tiny experiments, where process start-up dominates.
\end{systemslens}

% =====================================================================
\section{A typical training loop}\label{sec:pytorch:loop}

With a model, a loss and a loader, the training loop has the same skeleton in every chapter of this book
(\cref{fig:pytorch:loop}):

\begin{lstlisting}[caption={The canonical training loop.},label={lst:pytorch:loop}]
model = MLP(in_dim=2, out_dim=2)
optimizer = torch.optim.SGD(model.parameters(), lr=0.5)
for epoch in range(3):
    model.train()                          # enable dropout etc.
    for xb, yb in train_loader:
        logits = model(xb)                 # 1. forward
        loss = F.cross_entropy(logits, yb) # 2. loss
        optimizer.zero_grad()              # 3. clear old gradients
        loss.backward()                    # 4. backward: fill .grad
        optimizer.step()                   # 5. update parameters
    model.eval()                           # disable dropout for evaluation
    print(epoch, compute_accuracy(model, test_loader))

@torch.no_grad()
def compute_accuracy(model, loader):
    correct = total = 0
    for xb, yb in loader:
        correct += (model(xb).argmax(dim=-1) == yb).sum().item()
        total += yb.numel()
    return correct / total
\end{lstlisting}

\begin{figure}[t]
\centering
\begin{tikzpicture}[font=\footnotesize,
  s/.style={draw,rounded corners=3pt,minimum width=2.2cm,minimum height=0.8cm,align=center},
  arr/.style={-{Stealth[length=2mm]},thick,gray!70}]
\node[s,fill=keygreenbg,draw=keygreen] (load) at (90:2.6) {next batch};
\node[s,fill=sysbluebg,draw=sysblue] (fwd) at (18:2.9) {forward\\logits};
\node[s,fill=sysbluebg,draw=sysblue] (loss) at (-54:2.6) {loss};
\node[s,fill=warnredbg,draw=warnred] (bwd) at (-126:2.6) {zero\_grad +\\backward};
\node[s,fill=napkinbg,draw=napkin] (step) at (162:2.9) {optimizer\\step};
\draw[arr] (load) to[bend left=18] (fwd);
\draw[arr] (fwd) to[bend left=18] (loss);
\draw[arr] (loss) to[bend left=18] (bwd);
\draw[arr] (bwd) to[bend left=18] (step);
\draw[arr] (step) to[bend left=18] (load);
\node[align=center,font=\scriptsize] at (0,0) {one iteration\\$=$ one step\\[2pt]\texttt{model.train()}};
\node[draw,dashed,rounded corners=2pt,align=center,font=\scriptsize,fill=white] at (5.6,-1.6)
  {every $k$ steps / epoch:\\\texttt{model.eval()},\\validation, checkpoint};
\end{tikzpicture}
\caption{The training loop. Each iteration loads a batch, runs forward, computes the loss, clears stale
gradients, runs backward, and steps the optimiser. Evaluation and checkpointing happen periodically with
\code{model.eval()} and \code{torch.no\_grad()}.}
\label{fig:pytorch:loop}
\end{figure}

A few remarks generalise to LLM training. \code{model.train()} and \code{model.eval()} switch the behaviour of layers
such as dropout; forgetting \code{eval()} during generation makes outputs random for the wrong reason.
\code{optimizer.zero\_grad()} can sit before \code{backward()} or after \code{step()}, but must happen once per
optimiser step; \code{zero\_grad(set\_to\_none=True)} (the default) frees the gradient tensors instead of writing
zeros. The learning rate is the most important hyperparameter, and in \cref{ch:pretrain} we add warmup, cosine
decay and gradient clipping to this loop. Evaluating on a held-out set separates fitting from generalisation;
for LLMs we will track validation loss and, later, task accuracy.

% =====================================================================
\section{Saving and loading models}\label{sec:pytorch:save}

A module's \code{state\_dict()} is an ordered mapping from parameter (and buffer) names to tensors. Saving and
restoring it is the standard way to checkpoint:

\begin{lstlisting}[caption={Saving and restoring weights (and, for resuming, optimiser state).},label={lst:pytorch:save}]
torch.save(model.state_dict(), "model.pth")
model = MLP(2, 2)                                  # same architecture
model.load_state_dict(torch.load("model.pth", weights_only=True))

# resuming training needs the optimiser state too (e.g. Adam moments)
torch.save({"model": model.state_dict(), "opt": optimizer.state_dict(),
            "step": step}, "ckpt.pth")
\end{lstlisting}

\code{torch.save} uses Python's pickle format, which can execute arbitrary code on load, so the flag
\code{weights\_only=True} should be used for files you did not create. The ecosystem has largely moved to
\textbf{safetensors} for distributing weights: a simple header plus raw tensor bytes, safe to load and
memory-mappable, so a 16\,GB checkpoint can be opened instantly and paged in lazily. Both the GPT-2 and Qwen3
loaders in \code{llmsys} read safetensors files.

\begin{napkinmath}[Checkpoint sizes]
A training checkpoint holds far more than the weights. For a 7B model trained with mixed-precision AdamW, the
fp32 master weights and two Adam moments total $7\times10^9\times12=84$\,GB, plus 14\,GB of bf16 weights if stored.
At 2\,GB/s of write bandwidth to a parallel file system, saving takes about 50\,s, so frequent synchronous
checkpointing at scale wastes real GPU time. Sharded and asynchronous checkpointing (\cref{ch:pretrain,ch:distributed})
exist for this reason.
\end{napkinmath}

% =====================================================================
\section{Training on GPUs}\label{sec:pytorch:gpu}

\subsection{Computations on GPU devices}
Every tensor lives on a device, and an operation runs on the device of its inputs. Moving a tensor is
explicit, and mixing devices in one operation is an error:

\begin{lstlisting}[caption={Device placement rules.},label={lst:pytorch:device}]
device = torch.device("cuda" if torch.cuda.is_available() else "cpu")
a = torch.tensor([1., 2., 3.]).to(device)
b = torch.tensor([4., 5., 6.]).to(device)
print(a + b)            # computed on the GPU; result lives there too
c = torch.tensor([1., 1., 1.])
# a + c  -> RuntimeError: tensors on different devices (cuda:0 and cpu)
model = model.to(device)  # moves all parameters and buffers in place
\end{lstlisting}

On a machine with several GPUs, \code{"cuda:0"}, \code{"cuda:1"}, \ldots{} name individual devices, and the
environment variable \code{CUDA\_VISIBLE\_DEVICES} restricts and renumbers what a process can see.

\subsection{Single-GPU training}
Moving the loop of \cref{lst:pytorch:loop} to a GPU needs three changes: move the model once, move each batch
inside the loop, and leave everything else alone. For a tiny MLP the GPU will not be faster (each kernel does
almost nothing and launch overhead dominates), but for anything matrix-heavy, including every LLM, the
speedups are one to two orders of magnitude over a CPU.

\begin{lstlisting}[caption={The training loop on a GPU.},label={lst:pytorch:gpuloop}]
model = MLP(2, 2).to(device)
optimizer = torch.optim.SGD(model.parameters(), lr=0.5)
for epoch in range(3):
    model.train()
    for xb, yb in train_loader:
        xb, yb = xb.to(device, non_blocking=True), yb.to(device, non_blocking=True)
        loss = F.cross_entropy(model(xb), yb)
        optimizer.zero_grad()
        loss.backward()
        optimizer.step()
\end{lstlisting}

\subsection{Training with multiple GPUs}\label{sec:pytorch:ddp}
The simplest way to use several GPUs is \emph{data parallelism} with PyTorch's
\code{DistributedDataParallel} (DDP). One process is launched per GPU; each holds a full replica of the model
and receives a different slice of every global batch. After the backward pass, gradients are averaged across
processes with an \emph{all-reduce} collective, so every replica applies the same update and the replicas stay
identical (\cref{fig:pytorch:ddp}). With $P$ GPUs, an epoch takes roughly $1/P$ of the time, minus
communication overhead.

\begin{figure}[t]
\centering
\begin{tikzpicture}[font=\footnotesize,
  g/.style={draw=sysblue,fill=sysbluebg,rounded corners=2pt,minimum width=2.4cm,minimum height=2.4cm,align=center},
  d/.style={draw=keygreen,fill=keygreenbg,minimum width=0.8cm,minimum height=0.5cm,font=\scriptsize},
  arr/.style={-{Stealth[length=1.8mm]},thick,gray!70}]
\node[draw=gray,fill=gray!8,minimum width=4.4cm,minimum height=0.6cm] (batch) at (3.6,3.6) {global batch};
\foreach \i in {0,1,2} {
  \node[g] (g\i) at (\i*3.6,0.5) {};
  \node[font=\scriptsize\bfseries] at (\i*3.6,1.45) {GPU \i\ (rank \i)};
  \node[draw,fill=white,font=\scriptsize,minimum width=1.8cm] at (\i*3.6,0.85) {model replica};
  \node[font=\scriptsize] at (\i*3.6,0.3) {fwd + bwd};
  \node[draw=warnred,fill=warnredbg,font=\scriptsize,minimum width=1.8cm] (gr\i) at (\i*3.6,-0.25) {local grads $g_\i$};
  \node[d] (s\i) at (\i*3.6,2.6) {shard \i};
  \draw[arr] (batch.south) -- (s\i.north);
  \draw[arr] (s\i) -- (g\i.north);
}
\draw[very thick,labpurple,{Stealth[length=2mm]}-{Stealth[length=2mm]}] (gr0.south) -- ++(0,-0.55) -| (gr2.south);
\draw[very thick,labpurple] (gr1.south) -- ++(0,-0.55);
\node[font=\scriptsize,labpurple,fill=white] at (3.6,-1.1) {all-reduce: every rank gets $\bar g=\tfrac13(g_0+g_1+g_2)$, then identical optimiser steps};
\end{tikzpicture}
\caption{Distributed data parallelism. A \code{DistributedSampler} gives each rank a disjoint shard; each rank
computes gradients on its shard; an all-reduce averages them; identical updates keep replicas identical.}
\label{fig:pytorch:ddp}
\end{figure}

The code changes are modest. A process group is initialised with a rendezvous address, the model is
wrapped in \code{DDP}, and the loader uses a \code{DistributedSampler}. The complete script is
\code{labs/lab02\_pytorch\_roofline/ddp\_demo.py}; its heart is:

\begin{lstlisting}[caption={The essential pieces of DDP training (excerpt of \texttt{ddp\_demo.py}).},label={lst:pytorch:ddpcode}]
import torch.distributed as dist
from torch.nn.parallel import DistributedDataParallel as DDP
from torch.utils.data.distributed import DistributedSampler

def setup(rank, world_size):
    os.environ.setdefault("MASTER_ADDR", "localhost")   # rendezvous host
    os.environ.setdefault("MASTER_PORT", "29512")       # any free port
    backend = "nccl" if torch.cuda.is_available() else "gloo"
    dist.init_process_group(backend, rank=rank, world_size=world_size)
    if backend == "nccl":
        torch.cuda.set_device(rank)

def worker(rank, world_size, epochs):
    setup(rank, world_size)
    torch.manual_seed(0)                       # same init on every rank
    model = build_mlp(2, [16, 16], 2).to(device)
    ddp = DDP(model, device_ids=[rank] if device.type == "cuda" else None)
    sampler = DistributedSampler(ds, num_replicas=world_size, rank=rank)
    loader = DataLoader(ds, batch_size=16, sampler=sampler, pin_memory=True)
    for epoch in range(epochs):
        sampler.set_epoch(epoch)               # different shuffle each epoch
        for x, y in loader:
            loss = F.cross_entropy(ddp(x.to(device)), y.to(device))
            opt.zero_grad(); loss.backward(); opt.step()
    dist.destroy_process_group()

torch.multiprocessing.spawn(worker, args=(world_size, 3), nprocs=world_size)
\end{lstlisting}

In practice, scripts are usually launched with \code{torchrun --nproc\_per\_node=N script.py}, which spawns the
processes and sets \code{RANK}, \code{LOCAL\_RANK}, \code{WORLD\_SIZE}, \code{MASTER\_ADDR} and \code{MASTER\_PORT}. DDP
requires real processes, so it does not work from inside a Jupyter notebook. The NCCL backend implements the
collectives on NVIDIA GPUs over NVLink and InfiniBand; the Gloo backend runs them over TCP and is what the
lab uses on CPU.

\begin{systemslens}[What DDP actually does during backward]
DDP does not wait for the full backward pass and then all-reduce everything. It registers autograd hooks
that fire as each parameter's gradient becomes ready, groups gradients into \emph{buckets} (25\,MB by default),
and launches an asynchronous all-reduce per bucket while backward continues on earlier layers. Communication
thus overlaps with computation. A ring all-reduce moves about $2(P-1)/P$ times the gradient size per GPU, so
for a 1B-parameter model with bf16 gradients each step moves about 4\,GB per GPU; over NVLink at hundreds of
GB/s this hides easily behind compute, over 25\,Gb/s Ethernet it does not. DDP also assumes every replica
fits the full model, gradients and optimiser state on one GPU. \Cref{ch:distributed} removes that assumption
with ZeRO/FSDP, tensor and pipeline parallelism, and derives these communication costs properly.
\end{systemslens}

% =====================================================================
\section{The machine underneath: how a GPU is organised}\label{sec:pytorch:gpuarch}

We now switch from the API to the hardware. The details differ across generations and vendors, but the
organisation of an NVIDIA datacentre GPU is representative (\cref{fig:pytorch:hierarchy}); we use the H100
SXM as the running example.

\textbf{Streaming multiprocessors.} The chip consists of many identical \emph{streaming multiprocessors}
(SMs), 132 on an H100 SXM. Each SM runs many threads concurrently, organised in \emph{warps} of 32 threads that
execute the same instruction in lock-step (SIMT), and groups of warps form \emph{thread blocks} that are
scheduled onto one SM and can share fast on-chip memory. Each SM contains scalar/vector units and
\emph{tensor cores}, specialised units that compute small matrix products (for example $16\times8\times16$ tiles) in
one instruction. Nearly all of an H100's 989\,TFLOP/s of bf16 throughput comes from tensor cores; the
non-tensor fp32 throughput is more than an order of magnitude lower. A kernel that cannot be phrased as
matrix multiplication cannot use most of the chip's FLOPs.

\textbf{The memory hierarchy.} Data move through a hierarchy whose levels trade capacity for bandwidth and
latency: per-thread \emph{registers} (256\,KB per SM in total); per-SM \emph{shared memory}/L1 (up to about
228\,KB of configurable shared memory per SM on H100), explicitly managed by the kernel; a chip-wide \emph{L2
cache} (50\,MB); and off-chip \emph{high-bandwidth memory} (HBM: 80\,GB at 3.35\,TB/s). Host memory sits beyond
PCIe (tens of GB/s), and other GPUs are reachable over NVLink (900\,GB/s aggregate per H100).

\begin{figure}[t]
\centering
\begin{tikzpicture}[font=\footnotesize]
\def\lvl#1#2#3#4#5{% y, half-width, colour, label, right note
  \fill[#3!25,draw=#3] (-#2,#1) -- (#2,#1) -- (#2-0.55,#1+0.8) -- (-#2+0.55,#1+0.8) -- cycle;
  \node at (0,#1+0.4) {#4};
  \node[anchor=west,font=\scriptsize,align=left] at (4.1,#1+0.4) {#5};}
\lvl{3.2}{1.6}{warnred}{registers}{$\sim$256\,KB/SM $\cdot$ fastest, private to a thread}
\lvl{2.4}{2.15}{keygreen}{shared memory / L1}{$\le$228\,KB/SM $\cdot$ programmer-managed tiles}
\lvl{1.6}{2.7}{labpurple}{L2 cache}{50\,MB $\cdot$ shared by all SMs}
\lvl{0.8}{3.25}{napkin}{HBM}{80\,GB $\cdot$ 3.35\,TB/s}
\lvl{0.0}{3.8}{sysblue}{host DRAM / other GPUs}{PCIe $\sim$64\,GB/s $\cdot$ NVLink 900\,GB/s}
\draw[-{Stealth[length=2mm]},thick] (-4.3,0.2) -- (-4.3,3.9) node[midway,left,align=right,font=\scriptsize] {faster,\\smaller};
\end{tikzpicture}
\caption{The memory hierarchy of an H100 SXM (approximate published figures). Every kernel's performance is
determined by how often it must go down this pyramid. Fast kernels such as FlashAttention
(\cref{ch:attention}) restructure computation so that data loaded from HBM are reused many times in
shared memory and registers.}
\label{fig:pytorch:hierarchy}
\end{figure}

\textbf{Kernels.} A GPU kernel is a function executed by a grid of thread blocks. A typical PyTorch
operation, say \code{torch.add(x, y)}, launches one kernel that reads $x$ and $y$ from HBM, adds them in
registers, and writes the result back to HBM. A matrix multiplication launches a cuBLAS kernel that tiles the
output, stages input tiles in shared memory, and feeds tensor cores from registers. The difference in how much
work each loaded byte supports is the subject of the next section.

% =====================================================================
\section{The roofline model}\label{sec:pytorch:roofline}

A kernel needs time for its arithmetic and time to move its data. If the two overlap perfectly (a good
first approximation for well-written kernels), its runtime is bounded below by the larger:
\begin{equation}
t \;\ge\; \max\!\left(\frac{F}{\pi},\; \frac{Q}{\beta}\right),
\label{eq:pytorch:roofline}
\end{equation}
where $F$ is the number of FLOPs, $Q$ the bytes moved to and from HBM, $\pi$ the peak FLOP/s and $\beta$ the peak
bandwidth. Dividing $F$ by $t$ gives the attainable throughput as a function of the kernel's
\emph{arithmetic intensity} $I = F/Q$ (FLOPs per byte):
\begin{equation}
\text{attainable FLOP/s} \;=\; \min(\pi,\; I\cdot\beta).
\label{eq:pytorch:attain}
\end{equation}
Plotted on log--log axes, this is a sloped line (bandwidth-bound region) meeting a flat line (compute-bound
region) at the \emph{ridge point} $I^\star=\pi/\beta$: the \emph{roofline} (Williams, Waterman and Patterson,
2009). Kernels left of the ridge are memory-bound; making their arithmetic faster does nothing, only moving
fewer bytes helps. Kernels right of it are compute-bound.

\begin{figure}[t]
\centering
\begin{tikzpicture}
\begin{axis}[width=0.92\linewidth,height=7.2cm,xmode=log,ymode=log,
  xmin=0.05,xmax=5000,ymin=0.1,ymax=3000,
  xlabel={arithmetic intensity $I$ (FLOP/byte)},ylabel={attainable TFLOP/s},
  xlabel style={font=\footnotesize},ylabel style={font=\footnotesize},
  tick label style={font=\scriptsize},legend style={font=\scriptsize,at={(0.98,0.03)},anchor=south east},
  grid=major,grid style={gray!20}]
\addplot[very thick,sysblue,domain=0.05:5000,samples=400] {min(989,x*3.35)};
\addlegendentry{H100 SXM bf16 ($I^\star\approx295$)}
\addplot[very thick,napkin,dashed,domain=0.05:5000,samples=400] {min(312,x*2.039)};
\addlegendentry{A100 80GB bf16 ($I^\star\approx153$)}
\addplot[only marks,mark=*,mark size=2pt,warnred] coordinates {(0.1667,0.558) (1.0,3.35) (4096/3,989) (128/3,142.9) (4,13.4)};
\node[font=\scriptsize,anchor=north west] at (axis cs:0.19,0.5) {elementwise add (bf16), $I=\tfrac16$};
\node[font=\scriptsize,anchor=north west] at (axis cs:1.15,3.0) {decode GEMV, $I\approx1$};
\node[font=\scriptsize,anchor=north east] at (axis cs:2500,800) {GEMM $4096^3$, $I\approx1365$};
\node[font=\scriptsize,anchor=north west] at (axis cs:48,135) {GEMM $128^3$, $I\approx43$};
\node[font=\scriptsize,anchor=north west] at (axis cs:4.6,12) {softmax / norm, $I\sim4$};
\end{axis}
\end{tikzpicture}
\caption{Rooflines for H100 and A100 in bf16 (published peak numbers) with typical LLM kernels placed at
their arithmetic intensity (computed, not measured). Only large matrix multiplications reach the compute
roof; elementwise operations, normalisations and batch-1 decoding live on the bandwidth slope.}
\label{fig:pytorch:roofline}
\end{figure}

\subsection{Arithmetic intensity of the kernels that matter}
\paragraph{Matrix multiplication.} For $\mat{C}=\mat{A}\mat{B}$ with $\mat{A}\in\R^{m\times k}$, $\mat{B}\in\R^{k\times n}$
stored with $b$ bytes per element, $F=2mnk$ and, at best, $Q=b(mk+kn+mn)$ (each operand read once, the result
written once), so
\begin{equation}
I_{\text{GEMM}} = \frac{2mnk}{b(mk+kn+mn)}\;\xrightarrow{m=n=k}\;\frac{2n}{3b}.
\end{equation}
In bf16 a square GEMM has $I=n/3$: it crosses the H100 ridge near $n\approx885$. Large training GEMMs are
compute-bound; small ones are not.

\paragraph{Matrix--vector products.} With $m=1$ (one token's activations times a weight matrix),
$I=2nk/(b(k+kn+n))\approx 2/b$: about 1\,FLOP/byte in bf16. This is the operation at the heart of batch-1
decoding, and it sits at about $1/300$ of the H100's ridge point. With a batch of $B$ tokens sharing the same
weights, $m=B$ and $I\approx 2B/b$ for $B\ll k,n$; batching is the lever that slides decoding up the roofline.

\paragraph{Elementwise operations.} \code{y = x + z} in bf16 does 1 FLOP per element and moves 6 bytes:
$I=1/6$. Activation functions, residual additions, dropout, and the scaling in attention are all like this.
Normalisations and softmax perform a handful of FLOPs per element over two passes, $I$ of order 1--10. All are
deeply memory-bound, so the only way to speed them up is to avoid materialising their inputs and outputs in
HBM, by \emph{fusing} them into neighbouring kernels.

\begin{napkinmath}[Fusion pays for memory-bound chains]
Take a bf16 activation tensor of $8\times4096\times4096$ elements (256\,MiB) passing through bias-add, GELU and
dropout as three separate kernels: each reads and writes the tensor, about $3\times2\times256$\,MiB $\approx$ 1.6\,GB of
traffic, $\approx$0.48\,ms at 3.35\,TB/s. Fused into one kernel it reads once and writes once: 0.54\,GB,
$\approx$0.16\,ms, a $3\times$ speedup with identical FLOPs. This is exactly what \code{torch.compile}'s
code generator does automatically (\cref{sec:pytorch:async}) and what we do by hand in Triton in
\cref{app:triton}.
\end{napkinmath}

\begin{systemslens}[Reading the rest of the book through the roofline]
Many optimisations in later chapters are roofline arguments. FlashAttention (\cref{ch:attention}) raises
attention's intensity by never writing the $T\times T$ score matrix to HBM. The KV cache (\cref{ch:inference})
trades recomputation (FLOPs) for memory traffic. Grouped-query attention (\cref{ch:modern}) shrinks the
bytes read per decoded token. Batching and continuous batching (\cref{ch:serving}) raise decoding's $I$ by
sharing weight reads. Weight quantisation (\cref{ch:efficient}) cuts $Q$ directly. Speculative decoding
verifies several tokens per weight read. Each is a move to the right on \cref{fig:pytorch:roofline}, or a
reduction of $Q$ at fixed $F$.
\end{systemslens}

% =====================================================================
\section{The execution model: asynchrony, overhead and compilation}\label{sec:pytorch:async}

\subsection{Kernels are launched asynchronously}
When Python calls a CUDA operation, PyTorch does the bookkeeping on the CPU (argument checks, dispatch,
allocation), enqueues the kernel on a \emph{CUDA stream}, and returns immediately, before the GPU has done any
work. The CPU races ahead, enqueuing more kernels, while the GPU drains the queue in order. When the CPU is
ahead, launch overhead is hidden and the GPU is always busy. When individual kernels are tiny (small
tensors, many elementwise ops), the GPU finishes each kernel faster than the CPU can enqueue the next, and the
program becomes \emph{overhead-bound}: the GPU idles between kernels (\cref{fig:pytorch:async}).

\begin{figure}[t]
\centering
\begin{tikzpicture}[font=\scriptsize,x=0.62cm,y=0.55cm]
\node[anchor=east,font=\footnotesize\bfseries] at (-0.3,3) {CPU};
\node[anchor=east,font=\footnotesize\bfseries] at (-0.3,2) {GPU};
\node[anchor=west,font=\footnotesize] at (0,4.0) {(a) large kernels: CPU runs ahead, GPU saturated};
\foreach \i in {0,...,5} \fill[sysblue!50] (\i*0.9,2.75) rectangle ++(0.6,0.5);
\foreach \i in {0,...,3} \fill[napkin!70] (0.6+\i*2.6,1.75) rectangle ++(2.5,0.5);
\draw[-{Stealth},gray] (0,1.3) -- (11.5,1.3) node[right] {time};
\begin{scope}[yshift=-2.9cm]
\node[anchor=east,font=\footnotesize\bfseries] at (-0.3,3) {CPU};
\node[anchor=east,font=\footnotesize\bfseries] at (-0.3,2) {GPU};
\node[anchor=west,font=\footnotesize] at (0,4.0) {(b) tiny kernels: GPU waits for launches (overhead-bound)};
\foreach \i in {0,...,10} \fill[sysblue!50] (\i*1.0,2.75) rectangle ++(0.9,0.5);
\foreach \i in {0,...,10} \fill[napkin!70] (0.9+\i*1.0,1.75) rectangle ++(0.2,0.5);
\draw[-{Stealth},gray] (0,1.3) -- (11.5,1.3) node[right] {time};
\draw[decorate,decoration={brace,amplitude=3pt,mirror},thick] (1.1,1.65) -- (1.9,1.65) node[midway,below=3pt] {idle};
\end{scope}
\begin{scope}[yshift=-5.8cm]
\node[anchor=east,font=\footnotesize\bfseries] at (-0.3,3) {CPU};
\node[anchor=east,font=\footnotesize\bfseries] at (-0.3,2) {GPU};
\node[anchor=west,font=\footnotesize] at (0,4.0) {(c) same work fused or captured in a CUDA graph: one launch};
\fill[sysblue!50] (0,2.75) rectangle ++(0.9,0.5);
\fill[napkin!70] (0.9,1.75) rectangle ++(2.4,0.5);
\draw[-{Stealth},gray] (0,1.3) -- (11.5,1.3) node[right] {time};
\end{scope}
\node[anchor=west] at (0,-5.1) {\textcolor{sysblue!70}{\rule{8pt}{5pt}} launch + dispatch \quad \textcolor{napkin}{\rule{8pt}{5pt}} kernel execution};
\end{tikzpicture}
\caption{Asynchronous execution. (a) With large kernels, the CPU enqueues work faster than the GPU consumes it,
so launch overhead is hidden. (b) With tiny kernels, each launch costs more CPU time than the GPU spends
executing, leaving the GPU idle. (c) Fusion and CUDA graphs collapse many launches into one.}
\label{fig:pytorch:async}
\end{figure}

\begin{napkinmath}[When does launch overhead matter?]
A Python-level PyTorch op costs on the order of 5--20\,\textmu s of CPU time; a kernel that moves 1\,MB at
3.35\,TB/s runs in 0.3\,\textmu s. Decoding one token of Qwen3-0.6B touches 28 layers with a dozen or more small
kernels each: several hundred launches, i.e.\ milliseconds of CPU time, versus a bandwidth floor of about
0.36\,ms ($1.2\,\text{GB}/3.35\,\text{TB/s}$). Eager-mode decoding of small models is therefore often
CPU-bound by a large factor. This is why serving engines capture the decode step in a CUDA graph and why
\cref{ch:inference} shows large speedups from \code{torch.compile} on small models.
\end{napkinmath}

\subsection{Synchronisation points}
Because work is asynchronous, anything that needs a \emph{value} computed on the GPU forces the CPU to wait
for the queue to drain: \code{.item()}, \code{.tolist()}, \code{print(tensor)}, \code{.cpu()}, and Python control flow
that depends on tensor values (\code{if loss > 0:}). An innocent \code{loss.item()} for logging every step
serialises CPU and GPU and destroys overlap. Log every $k$ steps, or accumulate on the device. The same
asynchrony makes naive timing wrong: \code{time.time()} around a CUDA call measures enqueue time, not
execution time, unless the device is synchronised first (\cref{sec:pytorch:profiling}).

\subsection{Streams and overlap}
A stream is an in-order queue of kernels and copies. Work in different streams may overlap, which is how
communication (NCCL kernels on their own stream), host-to-device copies (with pinned memory) and compute run
concurrently. DDP's bucketed all-reduce and the loader's \code{non\_blocking} copies exploit this; in
\cref{ch:distributed,ch:serving} we schedule overlap explicitly.

\subsection{The caching allocator}
\code{cudaMalloc} and \code{cudaFree} are slow and synchronising, so PyTorch keeps freed blocks in a cache and
reuses them. Consequently \code{nvidia-smi} reports the cache's size rather than live tensors, and
\code{torch.cuda.memory\_allocated()} versus \code{torch.cuda.memory\_reserved()} tell the two apart. Allocation
patterns with many differently sized tensors (variable sequence lengths are the classic case) fragment the
cache, which can produce out-of-memory errors with plenty of nominally free memory. The KV-cache designs of
\cref{ch:inference,ch:serving} preallocate to avoid exactly this.

\subsection{Graphs and compilation}
Eager execution is flexible but pays per-op overhead and cannot optimise across operations. Two mechanisms
recover performance. \textbf{CUDA graphs} record a fixed sequence of kernel launches once and replay it with a
single launch; they require static shapes and static memory addresses, which is why they suit the
fixed-shape decode step of a serving engine. \textbf{\code{torch.compile}} captures Python-level programs into a
graph (TorchDynamo), derives the backward graph ahead of time (AOTAutograd), and generates fused kernels
(TorchInductor, emitting Triton on GPUs and C++/OpenMP on CPUs), optionally using CUDA graphs
(\code{mode="reduce-overhead"}). The first call is slow, since it compiles, and changing input shapes can trigger
recompilation. We use compilation for inference in \cref{ch:inference} and for training in \cref{ch:pretrain}.

\begin{lstlisting}[caption={Compiling a module. Same semantics, fewer and larger kernels.},label={lst:pytorch:compile}]
model = MLP(4096, 4096, hidden=(16384,)).to(device)
fast = torch.compile(model)                    # default: fuse ops, Triton kernels
x = torch.randn(64, 4096, device=device)
fast(x)                                        # first call compiles (seconds)
fast(x)                                        # later calls are fast
fast_rt = torch.compile(model, mode="reduce-overhead")  # + CUDA graphs
\end{lstlisting}

% =====================================================================
\section{Measuring and profiling}\label{sec:pytorch:profiling}

Performance claims without careful measurement are folklore. Three rules make timings trustworthy:
\textbf{warm up} (the first calls pay for allocator growth, cuBLAS heuristic selection, or compilation);
\textbf{synchronise} before reading the clock (asynchrony again); and \textbf{repeat and take the median}
(robust to noise from other processes). \code{llmsys.profiling.benchmark} implements exactly this:

\begin{lstlisting}[caption={Honest timing (from \texttt{llmsys/profiling.py}).},label={lst:pytorch:bench}]
def benchmark(fn, warmup=3, iters=10, min_time=0.0):
    for _ in range(warmup):
        fn()
    _sync()                                    # drain queued GPU work
    times, spent = [], 0.0
    while len(times) < iters or spent < min_time:
        t0 = time.perf_counter()
        fn()
        _sync()                                # wait for this call's kernels
        dt = time.perf_counter() - t0
        times.append(dt); spent += dt
    return statistics.median(times)
\end{lstlisting}

From a time, derive rates and compare them with the roofline. For a GEMM, achieved FLOP/s
$= 2mnk/t$ and achieved bandwidth $= b(mk+kn+mn)/t$:

\begin{lstlisting}[caption={Placing a measured kernel on the roofline.},label={lst:pytorch:place}]
from llmsys.profiling import benchmark, matmul_stats, Roofline
a = torch.randn(4096, 4096, device="cuda", dtype=torch.bfloat16)
b = torch.randn(4096, 4096, device="cuda", dtype=torch.bfloat16)
t = benchmark(lambda: a @ b)
s = matmul_stats(4096, 4096, 4096, t, bytes_per_el=2)
roof = Roofline(peak_flops=989e12, peak_bw=3.35e12)
print(f"{s['flops_per_s']/1e12:.0f} TFLOP/s of "
      f"{roof.attainable(s['intensity'])/1e12:.0f} attainable "
      f"({roof.bound(s['intensity'])}-bound)")
\end{lstlisting}

When a whole training step is slow, the question is \emph{which} kernels dominate. \code{torch.profiler} records
every operator and kernel with CPU and GPU times and can export a Chrome/Perfetto trace for timeline inspection;
\code{llmsys.profiling.top\_ops} returns the most expensive ops as a table:

\begin{lstlisting}[caption={Op-level profile of one training step.},label={lst:pytorch:profiler}]
from torch.profiler import profile, ProfilerActivity
acts = [ProfilerActivity.CPU, ProfilerActivity.CUDA]
with profile(activities=acts, record_shapes=True) as prof:
    loss = F.cross_entropy(model(x), y); loss.backward(); opt.step()
print(prof.key_averages().table(sort_by="cuda_time_total", row_limit=10))
prof.export_chrome_trace("step.json")          # open in ui.perfetto.dev
\end{lstlisting}

For kernel-level questions (achieved occupancy, memory throughput, tensor-core utilisation, stalls), NVIDIA's
Nsight Systems (timelines across CPU, GPU, NCCL and NVTX ranges) and Nsight Compute (per-kernel hardware
counters, including a roofline view) are the tools of record; \cref{app:triton} uses both.

\begin{pitfall}[Common measurement mistakes]
(1) Timing without synchronisation measures launch time. (2) Timing the first iteration measures
compilation and allocator warm-up. (3) Benchmarking with inputs that fit in L2 cache (50\,MB on H100)
overstates bandwidth; use buffers several times larger than L2 or flush it. (4) Reporting a mean over runs that
include a stray outlier: use medians and report spread. (5) Comparing dtypes without checking that the
fast path (tensor cores, tf32) was actually enabled. (6) On shared machines, noisy neighbours make CPU numbers
erratic; repeat until stable.
\end{pitfall}

\begin{systemslens}[Model FLOPs utilisation]
For an entire training run the headline metric is \emph{model FLOPs utilisation} (MFU): useful model FLOPs
per second (for a dense transformer, $6N$ per training token, ignoring recomputation) divided by the hardware
peak. If a 1.3B model trains at 150{,}000 tokens/s on 8 H100s, MFU $=150{,}000\times6\times1.3\times10^9/(8\times989\times10^{12})
\approx 15\%$, a sign of overheads (small kernels, communication, data loading) worth profiling. Well-tuned
large-scale runs report 35--55\%. MFU is the training analogue of the roofline fraction for a single kernel,
and \code{llmsys.flops.mfu} computes it.
\end{systemslens}

% =====================================================================
\section{Summary}\label{sec:pytorch:summary}
\begin{itemize}
\item PyTorch combines a tensor library, reverse-mode autograd, and deep-learning utilities. Models are
  \code{nn.Module}s; data flow through \code{Dataset}/\code{DataLoader}; training is the five-step loop of forward,
  loss, zero-grad, backward, step.
\item Tensors are views onto storage with strides; transposes and slices are free until a kernel needs
  contiguous memory. bf16 is the default LLM format: fp32 range, low precision.
\item Backward costs about twice the forward and requires saved activations, hence $6N$ FLOPs per training
  token and activation memory that scales with batch and sequence length.
\item DDP replicates the model per GPU, shards the data, and overlaps bucketed gradient all-reduces with
  backward.
\item GPUs deliver most FLOPs through tensor cores and feed them through a memory hierarchy. The roofline
  $\min(\pi, I\beta)$ says whether a kernel is limited by compute or bandwidth; most LLM operations other than
  large GEMMs are bandwidth-bound.
\item Execution is asynchronous: hidden launch overhead becomes visible with small kernels, and
  synchronisation points serialise CPU and GPU. Fusion, CUDA graphs and \code{torch.compile} reduce launches
  and memory traffic.
\item Measure with warm-up, synchronisation and medians; profile to find dominant kernels; summarise
  training efficiency with MFU.
\end{itemize}

% =====================================================================
\section{Exercises}\label{sec:pytorch:exercises}

\begin{exercise}[Strides]
For a contiguous tensor of shape $(2,3,4)$, write down its strides. What are the shape and strides after
\code{.transpose(0, 2)}? Explain why \code{.view(-1)} fails on the transposed tensor while
\code{.reshape(-1)} succeeds, and what the latter costs.
\end{exercise}

\begin{exercise}[bf16 rounding]
Compute the smallest update $\delta$ such that \code{torch.tensor(1.0, dtype=torch.bfloat16) + delta} differs
from 1.0. What does this imply for applying an update of size $\eta\, g = 10^{-4}$ to a weight of magnitude 1
stored in bf16? How do fp32 master weights fix this?
\end{exercise}

\begin{exercise}[Manual gradients]
Derive $\partial\mathcal{L}/\partial z = a-y$ for sigmoid followed by binary cross-entropy, and extend it to
a batch with the mean loss. Verify your formula against autograd numerically (this is Lab~2, Task~A1).
\end{exercise}

\begin{exercise}[Parameter counting]
Count the parameters of an MLP with layer sizes $784\to512\to256\to10$. How many bytes does it need for
weights, gradients and Adam states in fp32? How many FLOPs does one forward pass on a batch of 128 cost?
\end{exercise}

\begin{exercise}[Ridge points]
Compute the ridge point of an A100 (312\,TFLOP/s, 2.04\,TB/s) and an RTX~4090 (165\,TFLOP/s, 1.01\,TB/s) in bf16.
At what square GEMM size $n$ does each become compute-bound? Which GPU needs a larger decode batch to stop
being memory-bound?
\end{exercise}

\begin{exercise}[Intensity of a transformer MLP]
A GPT-style MLP computes $\GELU(\mat{X}\mat{W}_1)\mat{W}_2$ with $\mat{X}\in\R^{BT\times d}$,
$\mat{W}_1\in\R^{d\times4d}$, $\mat{W}_2\in\R^{4d\times d}$, $d=4096$, bf16. Compute FLOPs and minimal bytes for the
two GEMMs and the GELU as separate kernels for $BT=1$ and $BT=4096$. Which kernels are memory-bound in each
case? What is the saving from fusing GELU into the first GEMM's epilogue?
\end{exercise}

\begin{exercise}[DDP communication]
A 7B-parameter model is trained with DDP on 8 GPUs connected by NVLink (assume 300\,GB/s effective
all-reduce bandwidth per GPU) and, alternatively, on 8 nodes connected by 100\,Gb/s Ethernet. Estimate the
all-reduce time per step for bf16 gradients with a ring algorithm. If a step's compute takes 1.5\,s, which
setting can hide communication behind backward? (Also: would the model even fit for DDP?)
\end{exercise}

\begin{exercise}[Launch overhead]
Write a function applying 100 elementwise operations to a $1024$-element tensor and time it eagerly on a
GPU (or CPU) with proper synchronisation. Repeat with a $10^8$-element tensor. Then apply
\code{torch.compile}. Explain the three timings using \cref{fig:pytorch:async}.
\end{exercise}

\begin{exercise}[Profiling]
Profile one training step of the MLP of \cref{lst:pytorch:compile} with batch sizes 8 and 4096. Report the
top five ops by device time in each case and classify each as compute- or memory-bound using the roofline.
How does the fraction of time in \code{aten::mm}/\code{addmm} change?
\end{exercise}

\begin{exercise}[MFU sanity check]
A team reports training a 3B model at 2.1 million tokens/s on 64 H100s. Compute the implied MFU. Is the claim
plausible? What would you ask them about how ``tokens'' and ``FLOPs'' were counted?
\end{exercise}

% =====================================================================
\section{Lab 2: PyTorch essentials and an empirical roofline}\label{sec:pytorch:lab}

\begin{labbox}{2}{PyTorch essentials and an empirical roofline}
\labdir{lab02_pytorch_roofline}
\labsection{Goals}
Part~A makes the PyTorch workflow second nature: gradients by hand versus autograd, an \code{nn.Module},
a \code{DataLoader}, a training loop, checkpoints, and DDP on CPU processes. Part~B measures your own machine:
achieved FLOP/s and bandwidth, an empirical roofline, number formats, views versus copies, and an op-level
profile of a training step.

\labsection{Setup}
\begin{lstlisting}[style=shell]
cd code
pytest -q labs/lab02_pytorch_roofline                      # grades starter.py
LAB_IMPL=solution pytest -q labs/lab02_pytorch_roofline    # reference
\end{lstlisting}

\labsection{Tasks}
\textit{Part A: PyTorch.}
\begin{enumerate}
\item \ident{logistic_regression_grads(x, y, w, b)}: hand-derived gradients of mean BCE; the test compares
  them with autograd to $10^{-6}$.
\item \ident{build_mlp(in_dim, hidden, out_dim)}: an \code{nn.Sequential} returning logits; the test checks shapes
  and the exact parameter count.
\item \ident{train_classifier} and \ident{accuracy}: the loop of \cref{lst:pytorch:loop} with SGD; it must reach
  over 90\% test accuracy on the two-blob \code{ToyDataset}.
\item \ident{save_and_reload}: a \code{state\_dict} round trip that reproduces the accuracy exactly.
\item Run \code{ddp\_demo.py --world-size 2}. The provided test checks that the two replicas end with
  bit-identical weights. Explain why they do.
\end{enumerate}
\textit{Part B: the machine.}
\begin{enumerate}[resume]
\item \ident{matmul_intensity}, \ident{elementwise_intensity}, \ident{roofline_attainable}: the formulas of
  \cref{sec:pytorch:roofline}.
\item \ident{measure_matmul_sweep} and \ident{measure_bandwidth_sweep}: timed sweeps using
  \code{llmsys.profiling.benchmark}; \ident{empirical_roofline} picks the best observed FLOP/s and bytes/s.
\item \ident{shares_storage}: detect whether a tensor is a view, using storage pointers.
\item \ident{format_table}: max, epsilon and smallest normal of fp32/bf16/fp16 from \code{torch.finfo}.
\item \ident{profile_train_step}: profile forward, backward and step and return the top ops.
\end{enumerate}

\labsection{Measurements}
Run \code{python labs/lab02\_pytorch\_roofline/run.py} on your machine (\code{--device cuda --dtype bf16} on a
GPU) to produce \code{roofline.png}. Expect GEMM points to climb the bandwidth slope and flatten at the compute
roof for $n$ in the low thousands, and elementwise points to sit near $I=1/8$ (fp32) or $1/4$ (bf16) on the
bandwidth slope. On a GPU, compare your empirical peaks with the datasheet. Reaching 60--80\% of peak
bandwidth and 60--80\% of peak bf16 FLOP/s with plain PyTorch is typical. On a laptop CPU, expect noisy
results; repeat runs and note the variance.

\labsection{Deliverables}
Passing tests; your roofline plot annotated with the ridge point; a short table of the top five ops in a
training step at batch sizes 8 and 4096 with a one-line roofline explanation of each.

\labsection{Stretch goals}
(1) Add a bf16 and tf32 sweep and quantify the tensor-core speedup. (2) Time 100 tiny elementwise ops eagerly,
with \code{torch.compile}, and with \code{mode="reduce-overhead"}. (3) Use Nsight Compute's roofline section on
one GEMM and one elementwise kernel and compare with your empirical roofline. (4) Measure DDP scaling efficiency
for 1, 2 and 4 processes on a larger MLP and relate it to all-reduce volume.
\end{labbox}
